\section*{الفصل \ref{ch:g10:signals-and-waves} --- الإشارات والأمواج}

\begin{solution}{exo:g10:signals-and-waves:1}
$f = 1/T$: (أ) $1/\qty{0.020}{s} = \qty{50}{Hz}$؛
(ب) $1/\qty{4.0e-3}{s} = \qty{250}{Hz}$؛
(ج) $1/\qty{5.0e-5}{s} = \qty{20}{kHz}$.
\end{solution}

\begin{solution}{exo:g10:signals-and-waves:2}
$T = 1/f$: (أ) \qty{20}{ms}؛ (ب) $1/440 = \qty{2.3}{ms}$؛
(ج) $1/\qty{2.4e9}{Hz} = \qty{0.42}{ns}$.
\end{solution}

\begin{solution}{exo:g10:signals-and-waves:3}
\qty{12}{Hz}: \omterm{def:g10:signals-and-waves:ultrasound}{تحت صوتي}. و\qty{440}{Hz} و\qty{18}{kHz}: مسموع (والثاني
للآذان الفتيّة وحدها). و\qty{40}{kHz} و\qty{5}{MHz}: \omterm{def:g10:signals-and-waves:ultrasound}{فوق صوتي}.
\end{solution}

\begin{solution}{exo:g10:signals-and-waves:4}
$T = 5.0 \times \qty{2}{ms} = \qty{10}{ms}$، ومنه $f = \qty{100}{Hz}$؛
\omterm{def:g10:signals-and-waves:amplitude}{والسعة} $4.0 \times \qty{0.5}{V} = \qty{2.0}{V}$.
\end{solution}

\begin{solution}{exo:g10:signals-and-waves:5}
$d = 340 \times 6.0 = \qty{2040}{m} \approx \qty{2.0}{km}$. ويقطع الضوء
هذه المسافة في $2040/(3.00\times10^{8}) \approx \qty{7}{\micro s}$، أي أقلّ بمليون
مرة من \qty{6.0}{s}: فالومضة تعلّم لحظة الصاعقة.
\end{solution}

\begin{solution}{exo:g10:signals-and-waves:6}
ينزل \omterm{def:g10:signals-and-waves:sound}{الصوت} \emph{ويعود}، فيقطع $2d$:
$d = \frac{1500 \times 0.60}{2} = \qty{450}{m}$. وفوق الخندق:
$t = 2d/v = 2 \times 1200/1500 = \qty{1.6}{s}$.
\end{solution}

\begin{solution}{exo:g10:signals-and-waves:7}
$d_1 = \frac{3.00\times10^{8} \times 2.4\times10^{-4}}{2} =
\qty{36.0}{km}$ و$d_2 = \qty{35.7}{km}$: فقد قطعت الطائرة
\qty{300}{m} في \qty{1.0}{s}، ومنه $v = \qty{300}{m/s} \approx
\qty{1100}{km/h}$، وهي مقتربة.
\end{solution}

\begin{solution}{exo:g10:signals-and-waves:8}
$T = \qty{20}{mm}/\qty{25}{mm/s} = \qty{0.80}{s}$، ومنه
$f = \qty{1.25}{Hz}$: أي $75$ نبضة في الدقيقة. ومن أجل $100$ نبضة في الدقيقة،
$T = \qty{0.60}{s}$، أي\ $25 \times 0.60 = \qty{15}{mm}$: فيظهر تسرّع القلب
دون \qty{15}{mm}.
\end{solution}

\begin{solution}{exo:g10:signals-and-waves:9}
$T = 1/200 = \qty{5.0}{ms}$؛ وتدوم دورتان \qty{10}{ms}، ممتدتين على
$10$ تقسيمة: \omterm{def:g10:signals-and-waves:oscilloscope}{فقاعدة الزمن} \qty{1}{ms} لكل تقسيمة. ويجب أن تسع القمة
$4$ تقسيمات: \omterm{def:g10:signals-and-waves:oscilloscope}{فالكسب} $3.0/4 = \qty{0.75}{V}$ لكل تقسيمة على الأقل ---
خذ \qty{1}{V} لكل تقسيمة (فترتفع القمة $3.0$ تقسيمات).
\end{solution}

\begin{solution}{exo:g10:signals-and-waves:10}
$d = v\,t/2$: الجدار الأمامي
$\frac{1540 \times 4.0\times10^{-5}}{2} = \qty{3.1}{cm}$؛ والجدار الخلفي
$\frac{1540 \times 6.0\times10^{-5}}{2} = \qty{4.6}{cm}$؛ والسُّمك
نحو \qty{1.5}{cm}.
\end{solution}

\begin{solution}{exo:g10:signals-and-waves:11}
\omterm{def:g10:signals-and-waves:sound}{الصوت}: $30/340 = \qty{0.088}{s}$. والراديو:
$3.0\times10^{6}/(3.00\times10^{8}) = \qty{0.010}{s}$. فمستمع الراديو البعيد
يسمع الطبل أولًا، بنحو \qty{78}{ms}.
\end{solution}

\begin{solution}{exo:g10:signals-and-waves:12}
ما دام الخفاش يصدر \qty{3.0}{ms}، فهو أصمّ عن الأصداء، أي\ عن
كل ما هو ضمن $d = \frac{340 \times 3.0\times10^{-3}}{2} =
\qty{0.51}{m}$. والفراشة على \qty{1.7}{m}:
$t = 2 \times 1.7/340 = \qty{10}{ms}$. وكلما اقترب عاد \omterm{rem:g10:signals-and-waves:honest}{الصدى}
أبكر؛ ويجب أن تنتهي النبضة قبل وصوله، فوجب أن تقصر.
\end{solution}

\begin{solution}{exo:g10:signals-and-waves:13}
$d = \frac{340 \times 1.5}{2} = \qty{255}{m}$. وعند \qty{155}{m}:
$t = 2 \times 155/340 = \qty{0.91}{s}$. ويندمج \omterm{rem:g10:signals-and-waves:honest}{الصدى} في التصفيقة حين
$t < \qty{0.1}{s}$، أي\ ضمن
$d < \frac{340 \times 0.1}{2} = \qty{17}{m}$.
\end{solution}

\begin{solution}{exo:g10:signals-and-waves:14}
\emph{1.} $T = 3.0 \times \qty{5}{ms} = \qty{15}{ms}$،
$f = 1/0.015 \approx \qty{67}{Hz}$.
\emph{2.} $\qty{15}{ms}/\qty{2}{ms} = 7.5$ تقسيمة.
\emph{3.} تعرض الشاشة \qty{50}{ms}، ثم \qty{20}{ms}: أي $3$ دورة
كاملة ($50/15 = 3.3$)، ثم $1$ ($20/15 = 1.3$).
\end{solution}

\begin{solution}{exo:g10:signals-and-waves:15}
\emph{1.} $t = 2 \times 1.2\times10^{6}/(3.00\times10^{8}) =
\qty{8.0}{ms}$.
\emph{2.} $2 \times 1.2\times10^{6}/(2.0\times10^{8}) = \qty{12}{ms}$.
\emph{3.} لا يمتد الليف مستقيمًا بين الآلتين، وكل
موجّه وخادوم في الطريق يضيف تأخير معالجة.
\end{solution}

\begin{solution}{pb:g10:signals-and-waves:1}
\textbf{1.} $d = 340 \times 9.0 = \qty{3060}{m} \approx \qty{3.1}{km}$.

\textbf{2.} $3060/(3.00\times10^{8}) \approx \qty{10}{\micro s}$، أي أقصر
بمليون مرة من \qty{9.0}{s}: فالومضة تعلّم لحظة
الصاعقة.

\textbf{3.} في \qty{3}{s} يقطع \omterm{def:g10:signals-and-waves:sound}{الصوت} $340 \times 3 = \qty{1020}{m}
\approx \qty{1}{km}$: فالثواني مقسومة على ثلاثة تعطي الكيلومترات.

\textbf{4.} $d = 340 \times 6.0 = \qty{2040}{m}$. فقد اقتربت العاصفة
\qty{1020}{m} في \qty{300}{s}: أي $v = \qty{3.4}{m/s} \approx
\qty{12}{km/h}$، متجهة نحو القارب.

\textbf{5.} يُسمع الطرف القريب بعد $2000/340 = \qty{5.9}{s}$، ويُسمع الطرف
البعيد بعد $5000/340 = \qty{14.7}{s}$: فتدوم القرقعة نحو
\qty{8.8}{s}.

\textbf{6.} $d = 1500 \times 0.90/2 = \qty{675}{m}$.

\textbf{7.} $d = 1500 \times 0.20/2 = \qty{150}{m}$.

\textbf{8.} عائقان: سرب سمك على $1500 \times 0.16/2 =
\qty{120}{m}$ والقاع على \qty{150}{m}. فيسبح السرب
\qty{30}{m} فوق القاع.

\textbf{9.} يجب أن يعود \omterm{rem:g10:signals-and-waves:honest}{الصدى} قبل النبضة التالية:
$d_{\max} = 1500 \times 0.50/2 = \qty{375}{m}$. أما من ماء أعمق فيصل
\omterm{rem:g10:signals-and-waves:honest}{الصدى} \emph{بعد} النبضة التالية فيُنسب إليها: فتعرض
الشاشة قاعًا كاذبًا أضحل بكثير.

\textbf{10.} في الهواء، $340 \times 0.60/2 = \qty{102}{m}$ بدل
\qty{450}{m}. فلا يصير التأخير مسافة إلا متى عرفت الحامل
وسرعته في الوسط المقطوع.

\textbf{11.} $f = 1/0.86 = \qty{1.16}{Hz}$، أي\ $1.16 \times 60
\approx 70$ نبضة في الدقيقة.

\textbf{12.} $25 \times 0.86 = \qty{21.5}{mm}$.

\textbf{13.} $T = 15/25 = \qty{0.60}{s}$: أي $100$ نبضة في الدقيقة.

\textbf{14.} $T = 60/50 = \qty{1.2}{s}$؛ والتباعد
$25 \times 1.2 = \qty{30}{mm}$.

\textbf{15.} لا: \omterm{def:g10:signals-and-waves:period}{فالدور} ينزاح من نبضة إلى أخرى مع الجهد
والتنفس والتوتر النفسي. والطبيب يقرأ \omterm{def:g10:signals-and-waves:period}{الدور} (أي المعدل)،
وانتظامه، وشكل كل دورة --- فالتشخيص في
\omterm{def:g10:signals-and-waves:signal}{الإشارة}.

\textbf{16.} الراديو: $5.0\times10^{4}/(3.00\times10^{8}) =
\qty{0.17}{ms}$ --- أي آنيّ بالنسبة إلى الحواسّ البشرية. \omterm{def:g10:signals-and-waves:sound}{والصوت}:
$50\,000/340 \approx \qty{147}{s}$، أي دقيقتان ونصف دقيقة.

\textbf{17.} $t = 2.02\times10^{7}/(3.00\times10^{8}) = \qty{67}{ms}$.

\textbf{18.} $c \times \qty{1.0}{\micro s} = \qty{300}{m}$. ومن أجل تحديد
بدقة \qty{3.0}{m}: $3.0/(3.00\times10^{8}) = \qty{10}{ns}$ --- ولهذا
تحمل توابع GPS ساعات ذرّية.

\textbf{19.} التوابع في الخلاء، \omterm{def:g10:signals-and-waves:sound}{والصوت} يحتاج وسطًا: فلا
\omterm{def:g10:signals-and-waves:signal}{إشارة} تغادر \omterm{def:g10:universal-gravitation:satellite}{التابع} أصلًا. (وحتى لو سلّمنا بوجود هواء طول
الطريق، فإن \qty{20200}{km} عند \qty{340}{m/s} تزيد على \qty{16}{h} --- أي
تحديد موضع متأخر ساعات.)

\textbf{20.} \qty{340}{m/s} (\omterm{def:g10:signals-and-waves:sound}{الصوت} في الهواء)، و\qty{1500}{m/s} (\omterm{def:g10:signals-and-waves:sound}{الصوت} في
ماء البحر)، و$c = \qty{3.00e8}{m/s}$ (الراديو والضوء)؛ والقانون هو
$d = v\,t$، منصَّفًا في حالة \omterm{rem:g10:signals-and-waves:honest}{الصدى}. والعادة: اعرف حاملك، واعرف
سرعته في الوسط، ووقّت التأخير --- فالتأخير مسافة
متنكّرة.
\end{solution}
